\Needspace{4\baselineskip}
\begin{Example}
\label{ex:integration-operator-fixed-lower-endpoint}
Let $X=C_b([0,L])$ and define the \textbf{integration operator} $\mathcal{I}:X\to X$ by
\begin{equation}\label{eq:integration-operator-fixed-lower-endpoint}
\mathcal{I}(f)(x):=\int_0^xf(s)\,\mathrm{d}s,\quad x\in[0,L].
\end{equation}
Then $\mathcal{N}(\mathcal{I})=\{0\}$, and $\mathcal{R}(\mathcal{I})=\{g\in C_b^1([0,L])\mid g(0)=0\}$. The operator is bounded; the reader is asked to calculate its norm (Problem~\ref{ex:integration-operator-supremum-norm}).

Note that we could also think of this operator as mapping $L^2(0,L)$ to $L^2(0,L)$. (Hölder\textquotesingle s inequality tells us that $\mathcal{I}(f)(x)$ is well defined for every $f\in L^2(0,L)$ and every $x\in(0,L)$.) Once again using Hölder\textquotesingle s inequality we see that the operator is bounded since\[\begin{aligned}
\|\mathcal{I}(f)\|_{L^2(0,L)}^2&=\int_0^L|\mathcal{I}(f)(x)|^2\,\mathrm{d}x\\
&=\int_0^L\left(\int_0^xf(s)\,\mathrm{d}s\right)^2\,\mathrm{d}x\\
&\leq\int_0^L\left\{\int_0^x\mathrm{d}s\right\}\left\{\int_0^xf(s)^2\,\mathrm{d}s\right\}\,\mathrm{d}x\\
&\leq\dfrac{L^2}{2}\int_0^Lf(s)^2\,\mathrm{d}s\\
&=\dfrac{L^2}{2}\|f\|_{L^2(0,L)}^2.
\end{aligned}\]Using Theorem~\ref{thm:distribution-primitive-existence} we see that $\mathcal{N}(\mathcal{I})=\{0\}$.
\end{Example}

\Needspace{4\baselineskip}
\begin{Example}
\label{ex:differentiation-operator-domain-and-unboundedness}
We now define a \textbf{differentiation operator}. Of course, we can simply define an operator from, say, $C_b^1([0,L])$ to $C_b([0,L])$ by
\begin{equation}\label{eq:differentiation-operator-pointwise-definition}
Du(x):=\dfrac{\mathrm{d}u}{\mathrm{d}x}(x).
\end{equation}
With these spaces as domain and range it is easy to show that the differentiation operator is bounded. However, such a proof depends on the fact that we have imposed a ``stronger'' norm on the domain than on the range. In many contexts it turns out that such a restriction is unsatisfactory. Thus, let us consider the differentiation operator as mapping $L^2(0,L)$ to $L^2(0,L)$ with the domain of the operator restricted to $\widetilde{\mathcal{D}}(D):=C_b^1([0,L])$ and range $C_b([0,L])$; i.e., we consider\[L^2(0,L)\supset C_b^1([0,L])\ni u\mapsto\dfrac{\mathrm{d}u}{\mathrm{d}x}\in C_b([0,L])\subset L^2(0,L).\]It is easy to see that this operator is unbounded. We simply take
\begin{equation}\label{eq:unbounded-differentiation-exponential-sequence}
u_n(x)=e^{-nx}
\end{equation}
and calculate
\begin{equation}\label{eq:unbounded-differentiation-norm-ratio}
\dfrac{\left\|\dfrac{\mathrm{d}u_n}{\mathrm{d}x}\right\|_{L^2(0,L)}}{\|u_n\|_{L^2(0,L)}}=n.
\end{equation}
Note that unboundedness is not a real problem in inverting the differentiation operator. (We are obviously going to use an integration operator to accomplish that.) The most important difficulty is that the operator defined above has a nontrivial nullspace: the set of all constant functions. To solve this we simply restrict the domain using a boundary condition; for instance we can consider the domain $\mathcal{D}(D):=\{u\in C_b^1([0,L])\mid u(0)=0\}$. With this definition the differentiation operator is injective and the integration operator $\mathcal{I}$ defined above is its inverse.

Another area of concern is that the domain of the operator we have chosen is not ``optimal.'' We have left some very obvious candidates out: for instance, piecewise $C^1$ functions. This issue is addressed in the section on closed operators below (Section~\ref{subsec:closed-operators}).
\end{Example}
